\chapter{Setup}
\label{ch:setup}

Everything in this report runs on the released NLA for Qwen2.5-7B-Instruct at layer~20. The target model,
the verbalizer and the reconstructor are used exactly as published, except in the fine-tuning experiments,
where a small adapter is added to the verbalizer. This chapter describes the data and the three measuring
instruments: a detector for names, a probe that says what the activation still holds, and the statistics.

\section{Data}

I chose 60 entities, ten of each of six types (Table~\ref{tab:entities}). Three types are proper nouns and
three are common nouns, so that the results do not depend on capital letters or on how well a named-entity
recogniser works.

\begin{table}[t]
\centering
\small
\begin{tabular}{@{}lp{0.78\linewidth}@{}}
\toprule
Type & Entities \\
\midrule
Country & France, Japan, Egypt, Brazil, Canada, India, Germany, Italy, Mexico, Russia \\
City & Paris, Tokyo, Cairo, Chicago, London, Berlin, Moscow, Sydney, Toronto, Rome \\
Person & Shakespeare, Napoleon, Einstein, Mozart, Darwin, Lincoln, Churchill, Gandhi, Picasso, Beethoven \\
Animal & dog, cat, horse, lion, tiger, elephant, wolf, eagle, rabbit, dolphin \\
Food & bread, cheese, rice, pasta, chocolate, pizza, butter, honey, potato, sugar \\
Instrument & piano, guitar, violin, cello, flute, trumpet, saxophone, clarinet, harp, banjo \\
\bottomrule
\end{tabular}
\caption{The 60 entities.}
\label{tab:entities}
\end{table}

For each entity I collected 100 passages from the November 2023 English Wikipedia dump, 6{,}000 in all. A
passage qualifies if the entity appears in it exactly once, none of the other nine entities of its type
appears, at least 50 tokens come before the entity, and 16 tokens come after it. Articles whose title
contains the entity are skipped, so that the passage is about something else and merely mentions the entity,
and each article contributes at most one passage per entity.

The passages are split per entity: 30 for testing, and of the remaining 70, 14 for validation and 56 for
fitting. That gives 1{,}800 test, 840 validation and 3{,}360 fitting passages. Test passages are never used
to fit a probe or to fine-tune anything.

\section{Activations}

For every passage I ran the target model and saved the layer-20 activation at six positions: the entity's
last token, and 1, 2, 4, 8 and 16 tokens after it. I call these \emph{distances}. Distance~0 is where the
model has just read the entity; distance~16 is roughly a sentence later.

Before trusting any number, I checked that my pipeline reproduces known reconstruction quality. I had an
earlier pilot run of the released models, written separately and run on different hardware, which gave a
mean cosine of 0.907 between original and rebuilt activations. Repeating that run's recipe, the pipeline
used here gives 0.908, and explains 75.0\% of the variance, in line with the roughly 75\% the authors report
for this model. Getting there required fixing an off-by-one in which layer is read, described in
Appendix~B. On Wikipedia passages the mean cosine is 0.893, a little lower.

\section{Explanations}

Generating explanations is the expensive step, so the main measurements use a fixed subset of the test set:
nine test passages per entity, chosen by a seeded shuffle before any explanation was generated, at distances
0, 1, 4 and 16. That is 540 passages and 2{,}160 activations. For each activation the verbalizer wrote four
explanations, sampled independently at temperature~1, for 8{,}640 explanations in total. Four samples per
activation is what later makes it possible to ask whether the samples agree.

\section{Name Detector}

Every result depends on deciding, for an explanation and its passage, which names are true and which are
wrong. The rules are:

\begin{itemize}
\item The explanation \textbf{names the true entity} if the entity, or one of a small fixed set of variants
(\emph{wolves} for wolf, \emph{puppy} for dog), appears as a whole word.
\item A \textbf{wrong name} is any other entity of the same type that appears in the explanation and does
\emph{not} appear in the passage. A same-type name that is in the passage is real content the explanation
got right, so it is not counted against it.
\end{itemize}

``Any other entity of the same type'' has to reach well beyond my ten. For countries I use every country in
GeoNames, for cities every city above 300{,}000 people, for people the spans a named-entity recogniser tags
as persons (spaCy's small English model~\cite{spacy}), and for animals, foods and instruments the
corresponding branches of WordNet~\cite{miller1995wordnet}, keeping a word only where that sense is its most
common one. The broad lists matter: about four in five of the wrong names found later are outside my list of
60 (Section~\ref{sec:wrongnames}).

The detector is imperfect, and I would rather say how than hide it. A hand check of 30 flagged explanations
found four kinds of false alarm: category words (\emph{bird}, \emph{mammal}), meals (\emph{dinner}), words
for people in general (\emph{humans}), and ``persons'' that were not proper nouns. The detector used
throughout drops these. What remains is that the recogniser sometimes misses or invents a person, ambiguous
words such as \emph{bear} and \emph{bass} are never detected, and adjectives such as \emph{French} never
count as a country. These errors affect every condition I compare in the same way, since the same detector
scores all of them, but they do mean the absolute rates carry more uncertainty than their intervals show.

\section{Probe Baseline}

A falling true-name rate could mean two things: the verbalizer is failing to report something the activation
holds, or the activation no longer holds it. To tell these apart I need an independent estimate of what the
activation holds.

I use a linear probe~\cite{alain2016probes}. For each type $t$ and distance $\delta$, a ten-way logistic
regression on the standardised activation $\tilde h$ predicts which of the ten entities of that type is in
the passage:
\begin{equation}
p(e \mid h) = \operatorname{softmax}\bigl(A_{t,\delta}\, \tilde h + b_{t,\delta}\bigr)_e .
\end{equation}
Its regularisation strength is chosen on the validation passages from nine values between $10^{-5}$ and
$10^{3}$, and the probe is then refit on fitting and validation passages together and scored on the test
passages. Chance is 10\%.

The probe is a floor on what the activation holds, not a ceiling: a better reader might find more. It also
has an easier job than the verbalizer, since it chooses among ten known answers while the verbalizer writes
free text. I return to this when comparing the two.

\section{Statistics}

Passages about the same entity are not independent, so every interval in this report is a 95\% bootstrap
interval that resamples the 60 entities, with 10{,}000 resamples~\cite{efron1979bootstrap}. When two
conditions are compared on the same passages, I bootstrap the per-entity difference, which is much tighter
than comparing two separate intervals. All sampling uses fixed seeds, and every result is labelled as coming
from \emph{natural} text (real passages) or from \emph{planted} concepts (vectors I added by hand). Chapters
\ref{ch:fading} to \ref{ch:agreement} use natural text only.
